% =====================================================================
% pertemuan-12.tex
% Pertemuan 12 - Parser Generator Flex/Bison dan Aksi Semantik
% Sub-CPMK 4.3, 4.4 - Praktikum 3
% =====================================================================
\documentclass[aspectratio=169,10pt,t]{beamer}
\usepackage{config/beamerthemeubb}

\UBBmeeting{12}

\begin{document}

\UBBtitleframe{12}{Parser Generator Flex/Bison dan Aksi Semantik}{Praktikum 3}

% ---------------------------------------------------------------------
\begin{frame}{Capaian Pertemuan}
  \textbf{CPMK induk:} CPMK 04 --- Mengimplementasikan parser sederhana
  untuk validasi kode.

  \vspace{3pt}
  \begin{block}{Sub-CPMK yang dinilai}
    \footnotesize
    \begin{itemize}
      \setlength{\itemsep}{2pt}
      \item \textbf{Sub-CPMK 4.3:} Mengimplementasikan parser subset
            bahasa C dengan Flex dan Bison beserta aksi semantik,
            dengan berkas kode dan keluaran program sebagai bukti.
      \item \textbf{Sub-CPMK 4.4:} Membuat AST dan validator kode
            untuk subset bahasa C, lalu menguji keluaran pesan
            kesalahan sintaksis yang informatif.
    \end{itemize}
  \end{block}

  \vspace{2pt}
  {\footnotesize Praktikum 3 menilai 4.2, 4.3, dan 4.4; artefak 4.2
  dari Pertemuan 11 ikut dikumpulkan.}
\end{frame}

% ---------------------------------------------------------------------
\begin{frame}{Agenda Pertemuan (150 menit)}
  \small
  \begin{tabular}{@{}p{1.4cm}p{12.3cm}@{}}
    \toprule
    \textbf{Menit} & \textbf{Kegiatan} \\
    \midrule
    20 & Briefing Praktikum 3 dan cakupan artefak \\
    25 & Ceramah: arsitektur parser generator dan Bison \\
    35 & Praktik: deklarasi, aturan, dan aksi semantik \\
    35 & Praktik: integrasi Flex--Bison dan AST \\
    25 & Pengujian, diagnostik, dan perbandingan \\
    10 & Laporan dan penutup \\
    \bottomrule
  \end{tabular}

  \vspace{6pt}
  {\scriptsize\color{ubbblue!75}Porsi praktik 95 menit. Setiap
  mahasiswa menyerahkan grammar Bison, AST, dan keluaran pengujian.}
\end{frame}

% ---------------------------------------------------------------------
\begin{frame}{Posisi Pertemuan pada Rantai Proyek}
  \UBBchain{3}

  \vspace{2pt}
  \begin{itemize}
    \item Pertemuan ini menyelesaikan simpul \textbf{parser} dengan
          alat produksi: Bison membangun tabel LALR(1) dan AST.
    \item AST menjadi masukan analisis semantik pada Pertemuan 13--15.
    \item Praktikum 3 menghasilkan Milestone M3.
  \end{itemize}
\end{frame}

% ---------------------------------------------------------------------
\begin{frame}{Arsitektur Parser Generator}
  \begin{center}
    \begin{tikzpicture}[
        nd/.style={draw=ubbblue!40,fill=ubbblue!7,rounded corners=3pt,
          minimum width=3.0cm,minimum height=0.62cm,font=\scriptsize,
          text=ubbblue,align=center},
        data/.style={nd,draw=ubbgold,fill=ubbgold!20},
        ar/.style={-{Stealth[length=1.6mm]},draw=ubbblue!55,thick}
      ]
      \node[nd] (g) {Grammar dan aksi};
      \node[nd,right=0.5cm of g] (gen) {Parser generator};
      \node[nd,right=0.5cm of gen] (code) {Kode parser dan tabel};
      \node[data,below=0.5cm of code] (ast) {AST};
      \draw[ar] (g) -- (gen);
      \draw[ar] (gen) -- (code);
      \draw[ar] (code) -- (ast);
    \end{tikzpicture}
  \end{center}

  \begin{itemize}
    \item Pengguna menulis aturan produksi dan aksi semantik.
    \item Generator membangun tabel dan fungsi parser.
    \item Deskripsi bahasa terpisah dari algoritma parsing.
    \item Alat buku ini: Bison untuk parser, Flex untuk lexer.
  \end{itemize}
\end{frame}

% ---------------------------------------------------------------------
\begin{frame}{Keuntungan dan Keterbatasan}
  \begin{columns}[T,onlytextwidth]
    \begin{column}{0.48\textwidth}
      \begin{block}{Keuntungan}
        \footnotesize
        \begin{itemize}
          \item mengurangi kode yang ditulis tangan;
          \item tabel dibangun sistematis;
          \item perubahan grammar mudah diterapkan;
          \item mendukung grammar lebih besar.
        \end{itemize}
      \end{block}
    \end{column}
    \begin{column}{0.48\textwidth}
      \begin{block}{Keterbatasan}
        \footnotesize
        \begin{itemize}
          \item bergantung pada alat;
          \item pesan kesalahan kadang kurang kontekstual;
          \item perlu belajar membaca laporan konflik;
          \item versi alat memengaruhi hasil.
        \end{itemize}
      \end{block}
    \end{column}
  \end{columns}

  \vspace{4pt}
  {\footnotesize\color{ubbblue!80}Bison/Yacc memakai grammar LALR(1);
  Flex menghasilkan \eng{yylex} yang dipanggil Bison.}
\end{frame}

% ---------------------------------------------------------------------
\begin{frame}{Struktur Berkas Bison}
  \begin{enumerate}
    \setlength{\itemsep}{3pt}
    \item \textbf{Deklarasi}: token, tipe, precedence, dan kode C.
    \item \textbf{Aturan}: grammar dan \eng{semantic action}.
    \item \textbf{Kode pengguna}: fungsi bantu dan
          \texttt{yyerror}.
  \end{enumerate}

  \vspace{4pt}
  \begin{block}{Bagian deklarasi}
    \footnotesize
    \texttt{\%union} mendefinisikan tipe nilai semantik;
    \texttt{\%token} mendeklarasikan token; \texttt{\%type}
    menghubungkan nonterminal dengan tipe; \texttt{\%left} dan
    \texttt{\%right} menetapkan precedence serta associativity.
  \end{block}
\end{frame}

% ---------------------------------------------------------------------
\begin{frame}[fragile]{Deklarasi \file{parser.y}}
  \begin{columns}[T,onlytextwidth]
    \begin{column}{0.58\textwidth}
      \lstinputlisting[language=Bison,firstline=24,lastline=37,
        firstnumber=24,
        caption={Deklarasi \texttt{\%union}, token, dan
        \texttt{\%type}},
        label={lst:p12-decl}]{../buku_ajar/code/subsetc/src/parser.y}
    \end{column}
    \begin{column}{0.38\textwidth}
      \begin{itemize}
        \item Token membawa nilai bertipe \texttt{Token}.
        \item Nonterminal membawa \texttt{AstNode *}.
        \item \texttt{\%type} menghubungkan keduanya.
        \item \texttt{\%left} menetapkan precedence.
      \end{itemize}
    \end{column}
  \end{columns}
\end{frame}

% ---------------------------------------------------------------------
\begin{frame}[fragile]{Aturan dan Aksi: \texttt{program}}
  \begin{columns}[T,onlytextwidth]
    \begin{column}{0.58\textwidth}
      \lstinputlisting[language=Bison,firstline=53,lastline=65,
        firstnumber=53,
        caption={Aturan \texttt{program} dan aksinya},
        label={lst:p12-program}]{../buku_ajar/code/subsetc/src/parser.y}
    \end{column}
    \begin{column}{0.38\textwidth}
      \begin{itemize}
        \item Aksi membuat simpul \texttt{AST\_PROGRAM}.
        \item Daftar sementara memakai \texttt{next}.
        \item Setiap anak dipindahkan ke \texttt{children}.
        \item \texttt{\$\$} adalah hasil; \texttt{\$1}, \texttt{\$2}
              adalah nilai simbol.
      \end{itemize}
    \end{column}
  \end{columns}
\end{frame}

% ---------------------------------------------------------------------
\begin{frame}[fragile]{Bentuk Berkas Flex}
  \begin{columns}[T,onlytextwidth]
    \begin{column}{0.54\textwidth}
      \begin{lstlisting}[language=Flex,basicstyle=\ttfamily\scriptsize]
%{
#include "parser.tab.h"
%}
DIGIT [0-9]
%%
{DIGIT}+  { return T_INT_LIT; }
"+"       { return T_PLUS; }
[ \t\n]+  { /* lewati */ }
.         { return T_ERROR; }
%%
int yywrap(void) { return 1; }
      \end{lstlisting}
    \end{column}
    \begin{column}{0.42\textwidth}
      \begin{itemize}
        \item Tiga bagian: definisi, aturan, kode pengguna.
        \item Bagian aturan menghasilkan \texttt{yylex()}.
        \item Aksi mengisi \texttt{yylval} bila token punya nilai.
      \end{itemize}
    \end{column}
  \end{columns}
\end{frame}

% ---------------------------------------------------------------------
\begin{frame}{Nilai Semantik dan Pembangunan AST}
  \begin{itemize}
    \item \texttt{\$\$}: nilai hasil nonterminal kiri.
    \item \texttt{\$1}, \texttt{\$2}, \ldots: nilai simbol ruas kanan.
    \item \texttt{\$<token>1}: akses anggota union tertentu.
    \item Aksi dijalankan saat aturan direduksi.
    \item Aksi ini adalah \eng{translation scheme} Juarna: urutan aksi
          mengikuti simbol ruas kanan.
  \end{itemize}

  \vspace{5pt}
  \begin{block}{Pola pembangunan AST}
    \footnotesize
    Buat simpul dengan \texttt{ast\_new}, lalu tambahkan anak dengan
    \texttt{ast\_add\_child}. Daftar \texttt{next} dipakai selama
    parsing; \texttt{children} menyimpan struktur pohon final.
  \end{block}
\end{frame}

% ---------------------------------------------------------------------
\begin{frame}[fragile]{Aksi untuk Assignment, Call, dan Block}
  \begin{columns}[T,onlytextwidth]
    \begin{column}{0.58\textwidth}
      \lstinputlisting[language=Bison,firstline=135,lastline=149,
        firstnumber=135,
        caption={Aksi \texttt{assign\_or\_call}},
        label={lst:p12-assign}]{../buku_ajar/code/subsetc/src/parser.y}
    \end{column}
    \begin{column}{0.38\textwidth}
      \begin{itemize}
        \item \texttt{= expr} membangun simpul
              \texttt{AST\_ASSIGN}.
        \item \texttt{( args )} membangun simpul
              \texttt{AST\_CALL}.
        \item Argumen dipindahkan dari \texttt{next} ke
              \texttt{children}.
        \item Pola sama dipakai untuk blok.
      \end{itemize}
    \end{column}
  \end{columns}
\end{frame}

% ---------------------------------------------------------------------
\begin{frame}{AST untuk \texttt{a + b * c}}
  \begin{columns}[T,onlytextwidth]
    \begin{column}{0.46\textwidth}
      \centering
      \begin{tikzpicture}[
          nd/.style={draw=ubbblue!40,fill=ubbblue!7,rounded corners=2pt,
            minimum width=0.95cm,minimum height=0.52cm,font=\scriptsize},
          ar/.style={-{Stealth[length=1.4mm]},draw=ubbblue!55,thick}
        ]
        \node[nd] (plus) at (0,0) {\(+\)};
        \node[nd] (a) at (-1.3,-0.9) {\code{a}};
        \node[nd] (mul) at (1.2,-0.9) {\(*\)};
        \node[nd] (b) at (0.5,-1.8) {\code{b}};
        \node[nd] (c) at (1.9,-1.8) {\code{c}};
        \draw[ar] (plus) -- (a);
        \draw[ar] (plus) -- (mul);
        \draw[ar] (mul) -- (b);
        \draw[ar] (mul) -- (c);
      \end{tikzpicture}
    \end{column}
    \begin{column}{0.50\textwidth}
      \begin{itemize}
        \item Aksi \texttt{term} membangun \(*\) untuk
              \(b*c\).
        \item Aksi \texttt{additive} membangun \(+\) dengan anak
              \code{a} dan simpul \(*\).
        \item AST mencerminkan precedence tanpa tanda kurung.
        \item Simpul menyimpan token dan lokasi untuk diagnostik.
      \end{itemize}
    \end{column}
  \end{columns}
\end{frame}

% ---------------------------------------------------------------------
\begin{frame}[fragile]{Lokasi dan Pesan Kesalahan}
  \begin{columns}[T,onlytextwidth]
    \begin{column}{0.52\textwidth}
      \begin{lstlisting}[language=C,numbers=none,
        caption={Fungsi \texttt{yyerror} dengan lokasi}]
void yyerror(const char *msg) {
    fprintf(stderr, "%d:%d: %s\n",
            yylloc.first_line,
            yylloc.first_column, msg);
}
      \end{lstlisting}
    \end{column}
    \begin{column}{0.44\textwidth}
      \begin{itemize}
        \item \texttt{\%locations} menyediakan \texttt{@1}, \texttt{@2},
              dan \texttt{@\$}.
        \item Lokasi disalin ke simpul AST.
        \item Pesan memuat lokasi, penyebab, dan token ditemukan.
        \item ``syntax error'' saja tidak memenuhi Sub-CPMK 4.4.
      \end{itemize}
    \end{column}
  \end{columns}
\end{frame}

% ---------------------------------------------------------------------
\begin{frame}{Komponen Pesan Kesalahan Sintaksis}
  \small
  \setlength{\tabcolsep}{4pt}
  \begin{tabular}{@{}p{3.2cm}p{10.2cm}@{}}
    \toprule
    \textbf{Komponen} & \textbf{Contoh} \\
    \midrule
    Lokasi & baris 7, kolom 5 \\
    Penyebab & diharapkan \texttt{;} setelah ekspresi \\
    Token ditemukan & \texttt{\}} \\
    Saran perbaikan & tambahkan \texttt{;} sebelum \texttt{\}} \\
    \bottomrule
  \end{tabular}

  \vspace{5pt}
  \begin{block}{Pengujian}
    \footnotesize
    Pesan diuji pada kasus valid dan tidak valid. Diagnostik yang baik
    membantu pengguna menemukan kesalahan tanpa menebak.
  \end{block}
\end{frame}

% ---------------------------------------------------------------------
\begin{frame}{Integrasi Flex--Bison}
  \begin{center}
    \begin{tikzpicture}[
        nd/.style={draw=ubbblue!40,fill=ubbblue!7,rounded corners=3pt,
          minimum width=2.8cm,minimum height=0.62cm,font=\scriptsize,
          text=ubbblue,align=center},
        ar/.style={-{Stealth[length=1.6mm]},draw=ubbblue!55,thick}
      ]
      \node[nd] (src) {Program sumber};
      \node[nd,right=0.45cm of src] (flex) {Flex \texttt{yylex}};
      \node[nd,right=0.45cm of flex] (bison) {Bison \texttt{yyparse}};
      \node[nd,below=0.5cm of bison] (ast) {AST};
      \draw[ar] (src) -- (flex);
      \draw[ar] (flex) -- node[above]{\scriptsize token} (bison);
      \draw[ar] (bison) -- (ast);
    \end{tikzpicture}
  \end{center}

  \begin{itemize}
    \item \texttt{yylex} mengisi \texttt{yylval} dan mengembalikan kode
          token Bison.
    \item Driver memetakan jenis token lexer ke kode Bison.
    \item \texttt{yyparse} memakai tabel ACTION/GOTO LALR(1).
  \end{itemize}
\end{frame}

% ---------------------------------------------------------------------
\begin{frame}{Langkah Build}
  \begin{block}{Perintah build}
    \scriptsize
    \texttt{flex -o build/lexer\_flex.c src/lexer\_flex.l}\\[2pt]
    \texttt{bison -d -o build/parser.tab.c src/parser.y}\\[2pt]
    \texttt{gcc -Iinclude -Ibuild build/lexer\_flex.c
    build/parser.tab.c \textbackslash}\\
    \texttt{\hspace*{1.1cm}src/token.c src/error.c src/ast.c
    src/main\_bison.c -o build/scc}
  \end{block}

  \begin{itemize}
    \item Flex menghasilkan lexer C.
    \item Bison menghasilkan parser dan header token.
    \item GCC menautkan semuanya menjadi executable
          \texttt{scc}.
    \item Catat versi Flex, Bison, dan GCC pada laporan.
  \end{itemize}
\end{frame}

% ---------------------------------------------------------------------
\begin{frame}[fragile]{Driver Parser}
  \begin{columns}[T,onlytextwidth]
    \begin{column}{0.56\textwidth}
      \begin{lstlisting}[language=C,numbers=none,
        caption={Driver yang memanggil \texttt{yyparse}}]
int main(int argc, char **argv) {
    if (argc != 2) {
        fprintf(stderr, "penggunaan: %s berkas.sc\n",
                argv[0]);
        return 2;
    }
    if (!lexer_open(argv[1])) {
        return 2;
    }
    int status = yyparse();
    lexer_close();
    return status;
}
      \end{lstlisting}
    \end{column}
    \begin{column}{0.40\textwidth}
      \begin{itemize}
        \item Driver adalah alat uji.
        \item Membuka berkas dan menyiapkan lexer global.
        \item Memanggil \texttt{yyparse}.
        \item Mengembalikan status parsing.
        \item Pada milestone berikutnya, driver diganti alur
              semantik.
      \end{itemize}
    \end{column}
  \end{columns}
\end{frame}

% ---------------------------------------------------------------------
\begin{frame}{Bison dibandingkan Parser Tangan}
  \scriptsize
  \setlength{\tabcolsep}{4pt}
  \begin{tabular}{@{}p{3.0cm}p{5.6cm}p{5.4cm}@{}}
    \toprule
    \textbf{Aspek} & \textbf{Bison} & \textbf{Parser tangan} \\
    \midrule
    Algoritma & LALR(1) & Recursive descent \\
    Bentuk sumber & Berkas \texttt{.y} & Kode C langsung \\
    Aksi semantik & \texttt{\$\$}, \texttt{\$1} &
      Parameter dan nilai kembali \\
    Konflik & Dilaporkan alat & Dikelola manual \\
    Diagnostik & Perlu kode tambahan & Penuh, ditulis manual \\
    Integrasi lexer & Flex & Pemanggilan fungsi manual \\
    Cocok untuk & Grammar besar, proyek tim &
      Bahasa kecil, pembelajaran \\
    \bottomrule
  \end{tabular}

  \vspace{4pt}
  {\scriptsize\color{ubbblue!75}Keduanya dapat dipakai bersama: Flex
  menangani token, Bison menangani struktur.}
\end{frame}

% ---------------------------------------------------------------------
\begin{frame}{Pengujian Praktikum 3}
  \begin{columns}[T,onlytextwidth]
    \begin{column}{0.48\textwidth}
      \begin{block}{Kasus valid}
        \footnotesize
        \begin{itemize}
          \item deklarasi dan ekspresi;
          \item if/else dan while;
          \item panggilan fungsi;
          \item ekspresi campuran dengan kurung.
        \end{itemize}
      \end{block}
    \end{column}
    \begin{column}{0.48\textwidth}
      \begin{block}{Kasus tidak valid}
        \footnotesize
        \begin{itemize}
          \item titik koma hilang;
          \item kurung tidak seimbang;
          \item ekspresi tidak lengkap;
          \item deklarasi tanpa nama.
        \end{itemize}
      \end{block}
    \end{column}
  \end{columns}

  \vspace{4pt}
  {\footnotesize\color{ubbblue!80}Minimal 20 kasus uji. Untuk setiap
  kasus tidak valid, periksa lokasi dan isi pesan.}
\end{frame}

% ---------------------------------------------------------------------
\begin{frame}{Kegiatan Praktikum: M3}
  \begin{block}{Langkah kerja}
    \footnotesize
    \begin{enumerate}
      \setlength{\itemsep}{1pt}
      \item Lengkapi \file{src/parser.y}.
      \item Definisikan \texttt{\%union}, \texttt{\%token}, dan
            \texttt{\%type}.
      \item Tulis \eng{semantic action} untuk setiap aturan.
      \item Integrasikan \texttt{yylex} dari lexer.
      \item Jalankan Bison dan perbaiki konflik.
      \item Uji dan bandingkan dengan parser tangan.
    \end{enumerate}
  \end{block}

  \vspace{2pt}
  {\footnotesize\color{ubbblue!80}Kumpulkan tabel ACTION/GOTO SLR/LALR
  (4.2) dari Pertemuan 11 dan AST (4.4) yang dihasilkan hari ini.}
\end{frame}

% ---------------------------------------------------------------------
\begin{frame}{Rubrik Praktikum 3}
  \small
  \setlength{\tabcolsep}{4pt}
  \begin{tabular}{@{}p{3.4cm}p{8.6cm}p{1.8cm}@{}}
    \toprule
    \textbf{Aspek} & \textbf{Indikator} & \textbf{Bobot} \\
    \midrule
    Kebenaran & Grammar dan AST benar; konflik terselesaikan & 40\% \\
    Kualitas kode & Aksi ringkas, penamaan, komentar, tanpa
      peringatan & 25\% \\
    Kelengkapan & Ekspresi, pernyataan, deklarasi, dan AST & 20\% \\
    Pengujian & Kasus normal, batas, dan kesalahan terdokumentasi &
      15\% \\
    \bottomrule
  \end{tabular}

  \vspace{5pt}
  {\footnotesize\color{ubbblue!80}Rubrik ini sama dengan rubrik artefak
  praktikum pada bagian 10 RPS.}
\end{frame}

% ---------------------------------------------------------------------
\begin{frame}{Tugas dan Artefak}
  \begin{block}{Artefak Sub-CPMK 4.3 dan 4.4}
    Parser Flex/Bison dengan aksi semantik, berkas kode, keluaran
    program, AST, dan pesan kesalahan sintaksis yang informatif.
  \end{block}

  \small
  \begin{itemize}
    \item Grammar Bison tanpa konflik tak terduga.
    \item AST terbentuk untuk program valid.
    \item Minimal 20 kasus uji valid dan tidak valid.
    \item Diagnostik menyebut lokasi dan token yang diharapkan.
    \item Tabel ACTION/GOTO (4.2) dari Pertemuan 11 ikut dikumpulkan.
    \item Perbandingan dengan parser tangan terdokumentasi.
    \item Laporan singkat dan versi alat diserahkan.
  \end{itemize}

  \vspace{3pt}
  {\footnotesize\color{ubbblue!75}Praktikum 3 berbobot 12\% dan
  dilaporkan setelah Milestone M3 selesai.}
\end{frame}

% ---------------------------------------------------------------------
\begin{frame}{Rangkuman dan Jembatan}
  \begin{itemize}
    \item Parser generator memisahkan deskripsi grammar dari algoritma
          parsing.
    \item Bison memakai deklarasi token, tipe, precedence, aturan, dan
          aksi semantik.
    \item \texttt{\$\$} dan \texttt{\$1} menghubungkan nilai simbol
          untuk membangun AST.
    \item Flex menyediakan \texttt{yylex}; Bison menyediakan
          \texttt{yyparse}; keduanya ditautkan menjadi satu
          executable.
    \item Pesan kesalahan harus menyebut lokasi, penyebab, dan token
          yang ditemukan.
  \end{itemize}

  \vspace{4pt}
  \begin{exampleblock}{Pertemuan berikutnya}
    Pertemuan 13 beralih ke analisis semantik: attribute grammar,
    \eng{syntax-directed definition}, dan \eng{translation scheme}.
    AST dari Bison menjadi masukan aturan semantik.
  \end{exampleblock}
\end{frame}

\end{document}
